\section{How a large language model works}
\label{ch:llm}

A large language model (LLM) is a function that takes a sequence of tokens
and returns a probability distribution over the next token. Everything
else, chat, translation, reasoning, is that function applied repeatedly.
Figure~\ref{fig:pipeline} is the whole pipeline; the rest of the chapter
walks through each box, and says at each step what LFS reads.

\begin{figure}[h]\centering\includegraphics[width=\linewidth]{fig_llm_pipeline.png}
\caption{Text in, next-token probabilities out. LFS is computed from the hidden states in the middle and never asks the model to generate.}\label{fig:pipeline}\end{figure}

\subsection{Tokenization and fertility}

Text is split into \emph{tokens}: pieces of words drawn from a fixed
vocabulary of 30{,}000 to 250{,}000 entries, learned by byte-pair encoding
(BPE) or a similar scheme so that frequent strings become single tokens.
``The'' is one token; ``schläft'' may be two; an Arabic or Persian word may be
five or six because the vocabulary saw little of those scripts.

\emph{Fertility} is the number of tokens per word (or per character),
usually reported relative to English. High fertility means a language pays
more tokens for the same meaning, gets less content into a fixed context
window, and costs more per sentence. Fertility matters for LFS in two ways.
Mean pooling averages a different number of tokens per language, and the
tokenizer itself carries language identity into the embedding layer. That is
why fertility is a control variable in every validity analysis, and why the
paper reports that per-language contributions to the language variance
correlate with fertility at $+0.42$ to $+0.64$.

\subsection{Embeddings and the residual stream}

Each token id selects a row of the embedding matrix $E \in \R^{V\times D}$.
The resulting $D$-dimensional vector is the starting value of the
\emph{residual stream} for that position. Modern models add position
information through rotary embeddings inside attention rather than by adding
a position vector, but the picture is the same: position $t$ has a vector
$\mathbf x^{(0)}_t \in \R^D$.

The residual stream is a running sum. Each of the $n$ blocks reads the
stream, computes an update, and adds it back:
\[
\mathbf x^{(k)} = \mathbf x^{(k-1)} + \mathrm{Attn}_k\!\big(\mathrm{norm}(\mathbf x^{(k-1)})\big) + \mathrm{MLP}_k\!\big(\mathrm{norm}(\cdot)\big).
\]
The \emph{hidden state at layer $k$} is $\mathbf x^{(k)}$. Hugging Face
returns $n+1$ hidden states per token: index 0 is the embedding output and
index $k$ is the stream after block $k$. In the project's results a 28-block
model therefore has 29 ``layers'', and layer 0 is the embedding layer. This
additive structure is why it is meaningful to compare states across layers at
all: they all live in the same $D$-dimensional space and each layer is an
increment.

\begin{figure}[h]\centering\includegraphics[width=0.9\linewidth]{fig_llm_block.png}
\caption{One transformer block. The horizontal line is the residual stream; the two branches add their outputs back into it.}\label{fig:block}\end{figure}

\subsection{Attention}

Attention is how a position pulls information from earlier positions. For
each position compute a query $\mathbf q$, a key $\mathbf k$, and a value
$\mathbf v$ (three matrix multiplications). The score between position $i$
and an earlier position $j$ is $\mathbf q_i\cdot\mathbf k_j/\sqrt{d}$; a
softmax over $j$ turns the scores into weights (Figure~\ref{fig:attn}); the
output is the weighted sum of the values $\sum_j w_{ij}\mathbf v_j$. A
\emph{causal mask} forbids looking at later positions, which is what makes
the model a left-to-right predictor. Multi-head attention runs this several
times in parallel with smaller $d$ and concatenates.

\begin{figure}[h]\centering\includegraphics[width=0.62\linewidth]{fig_llm_attention.png}
\caption{Causal attention weights for one head: each row sums to one and only looks left.}\label{fig:attn}\end{figure}

The content-transfer experiment in the paper is attention in action: a
context sentence in Russian lowers the loss on an English target only if
attention at some layer moves information from the Russian tokens into the
English positions. The attention analysis measured exactly that mass, and
found it peaks just below the LFS minimum.

\begin{exercise}{single-head causal attention in numpy}
\begin{lstlisting}
T, D, d = 6, 32, 16
X = rng.normal(size=(T, D))                     # residual stream for 6 positions
Wq, Wk, Wv = (rng.normal(0, D**-0.5, (D, d)) for _ in range(3))
Q, K, V = X @ Wq, X @ Wk, X @ Wv
scores = Q @ K.T / np.sqrt(d)
scores[np.triu_indices(T, 1)] = -np.inf         # causal mask: no looking ahead
W = np.exp(scores - scores.max(1, keepdims=True)); W /= W.sum(1, keepdims=True)
out = W @ V                                     # (T, d): what each position gathered
print(np.round(W, 2))                           # lower-triangular, rows sum to 1
\end{lstlisting}
Now apply the same $W$ to \texttt{V} scaled by 100 and confirm \texttt{W} does
not change: attention weights depend on $Q$ and $K$, not on the scale of $V$.
\end{exercise}

\subsection{The MLP and normalization}

The MLP in each block is the two-matrix network of Chapter~\ref{ch:ml}, applied
to each position separately, with a hidden width of about $4D$ (SwiGLU
variants use a gated form). It is where much of a model's factual knowledge
appears to be stored, and it is the part that does not mix positions.

Normalization (LayerNorm or RMSNorm) rescales each position's vector to a
standard size before each branch. Modern models are \emph{pre-norm}: the
branch sees a normalized copy, but the residual stream itself is not
normalized, which is why a few coordinates of the stream can grow into
\emph{massive activations}, values hundreds of times larger than their
neighbors. Those coordinates dominate any unstandardized distance, overflow
fp16 sums, and are the reason LFS z-scores every coordinate jointly before
decomposing.

\subsection{From the last layer to a probability}

After the last block, a final norm and the \emph{unembedding} matrix $W_U \in
\R^{D\times V}$ turn the $D$-vector at each position into $V$ logits; softmax
gives the next-token distribution. Some models tie $W_U$ to the embedding
matrix. The \emph{logit lens} applies $W_U$ to an intermediate hidden state
and reads off the top token; it is a heuristic that presumes intermediate
states already live near the output space. The project measured that
presumption: mid-stack states reach a maximum cosine of about 0.10 to any
row of $W_U$ (0.20 at the final layer). Logit-lens readouts of middle layers,
including the ``LLMs think in English'' measure, rest on weak geometry.

\subsection{Training, base models, and post-training}

Training (Figure~\ref{fig:training}) is gradient descent on the next-token
cross-entropy over trillions of tokens, in batches, with Adam and a learning
rate that warms up and then decays. The model is always shown the true prefix
(\emph{teacher forcing}), so training is fully parallel across positions. The
output is a \emph{base} checkpoint: a next-token predictor, not a chatbot.
Post-training adds \emph{supervised fine-tuning} on instruction--response
pairs and often \emph{preference optimization} (RLHF, DPO) to produce an
\emph{instruct} model. Every model measured in the paper is a base checkpoint;
whether instruction tuning moves the LFS profile is an open, cheap
experiment.

\begin{figure}[h]\centering\includegraphics[width=\linewidth]{fig_llm_training.png}
\caption{The training loop. Nothing in it mentions languages; multilinguality is whatever the data induces.}\label{fig:training}\end{figure}

Two consequences for multilinguality. First, the data mix decides what is
learned: models trained mostly on English are English-centric by
construction, and the proposal's Aim 2 is about changing the data or the loss
to fix that. Second, ``scale'' (parameters and tokens) and ``recipe'' (mix,
tokenizer, curriculum) are confounded across labs, which is why the paper
calls the recipe-over-scale finding descriptive.

\subsection{Inference and scoring}

At inference the model produces a distribution at each step; you pick a
token (greedy, or sampling with a temperature) and repeat. For evaluation
you often do not generate at all: you feed a fixed continuation and read off
its probability. The teacher-forced NLL of a target sentence given a context
is exactly what the content-transfer test measures, and Belebele scoring
compares the log-likelihoods of the four answer options.

\subsection{Precision}

fp32 uses 32 bits per number; fp16 uses 16 with a small exponent range (max
about 65{,}000); bf16 uses 16 with fp32's exponent range but fewer digits.
Models are run in bf16 or fp16 for speed. Sums of squares of massive
activations overflow fp16, so pooling and every statistic in the project are
accumulated in fp32, and BLOOM, which was trained in bf16, is never run in
fp16. Two documented errors in the discrepancy ledger came from exactly
this; be ready to tell them.

\subsection{Interpretability handles you will use}

\begin{description}[leftmargin=2.6cm, style=nextline, itemsep=1pt]
\item[hidden states] the residual stream per layer, via \texttt{output\_hidden\_states=True}. The raw material of LFS.
\item[pooling] collapsing a sentence's token states into one vector: mean over tokens (used here), last token, or unit-normalized mean. Pooling is part of the measurement definition; absolute LFS values change with it, rankings much less.
\item[linear probe] a logistic regression trained on hidden states to predict a label (here: which language). High accuracy means the information is linearly present. Language stays decodable at 0.875 to 0.886 at every layer even where LFS is lowest.
\item[logit lens] unembed an intermediate state; see above for its limits.
\item[activation patching] copy a hidden state from one run into another and measure what changes; the causal tool Dumas et al.\ use to separate language from concept.
\item[CKA] a similarity index between two sets of representations, invariant to rotation and scale; used as a comparator in the fault test suite.
\end{description}

\begin{exercise}{hidden states for one sentence in three languages (run on CLSP)}
\begin{lstlisting}
# sbatch to a 2080Ti node; ~/venvs/mlm has transformers
import torch, numpy as np
from transformers import AutoTokenizer, AutoModel
name = "Qwen/Qwen3-0.6B-Base"
tok = AutoTokenizer.from_pretrained(name); model = AutoModel.from_pretrained(name, output_hidden_states=True).eval()
sents = {"eng": "The bat flew through the air.", "deu": "Die Fledermaus flog durch die Luft.", "fra": "La chauve-souris a volé dans les airs.", "spa": "El murciélago voló por el aire.", "rus": "<paste the Russian NTREX translation here>", "arb": "<paste the Arabic NTREX translation here>", "fas": "<paste the Persian NTREX translation here>"}
pooled = {}
for lang, s in sents.items():
    enc = tok(s, return_tensors="pt")
    with torch.no_grad(): hs = model(**enc).hidden_states       # tuple of (1, T, D), length n_layers+1
    pooled[lang] = np.stack([h[0].float().mean(0).numpy() for h in hs])  # (n_layers+1, D)
def cos(a, b): return a @ b / (np.linalg.norm(a) * np.linalg.norm(b))
for k in range(pooled["eng"].shape[0]):
    print(k, *[round(cos(pooled["eng"][k], pooled[l][k]), 3) for l in ("deu", "fra", "spa", "rus", "arb", "fas")])
\end{lstlisting}
You will see cosines that are high everywhere (anisotropy), which is why raw
cosine is a poor instrument and why the project standardizes and uses a
balanced grid instead of one sentence.
\end{exercise}

\subsection{The models in the paper, in one table}

\begin{center}\small
\begin{tabular}{llll}
\toprule
Family & Sizes measured & Training emphasis & Why it is in the grid \\
\midrule
Qwen3 & 0.6B, 1.7B, 4B, 8B & 119 languages, about 36T tokens & the within-family scaling series \\
OLMo-2 & 1B, 7B, 13B & mostly English, fully open data & English-centric control with known data \\
Mistral & 7B, Nemo 12B & English and European web & independent English-centric recipe \\
Llama-3.1 & 8B & English-heavy generalist & deepest dip in the grid \\
BLOOM & 1.7B, 7.1B & 46 languages, balanced, 2022 & the flat profile: coverage without convergence \\
Salamandra & 2B, 7B & 35 European languages, balanced, 2024 & modern balanced model; shallows with scale \\
EuroLLM & 1.7B & 35 European languages & targeted multilinguality \\
SmolLM2, Falcon3 & 360M, 1.7B; 3B, 7B & English-centric & the frozen held-out prospective test \\
granite, Yi & 8B; 9B & mixed & added under frozen predictions \\
Qwen2.5, XGLM, Sailor, Tower, Occiglot, mGPT & various & various & the 19-model expansion set \\
\bottomrule
\end{tabular}
\end{center}
